% Options for packages loaded elsewhere
\PassOptionsToPackage{unicode}{hyperref}
\PassOptionsToPackage{hyphens}{url}
\PassOptionsToPackage{dvipsnames,svgnames,x11names}{xcolor}
%
\documentclass[
  9pt,
  a4paper,
]{article}
\usepackage{amsmath,amssymb}
\usepackage{iftex}
\ifPDFTeX
  \usepackage[T1]{fontenc}
  \usepackage[utf8]{inputenc}
  \usepackage{textcomp} % provide euro and other symbols
\else % if luatex or xetex
  \usepackage{unicode-math} % this also loads fontspec
  \defaultfontfeatures{Scale=MatchLowercase}
  \defaultfontfeatures[\rmfamily]{Ligatures=TeX,Scale=1}
\fi
\usepackage{lmodern}
\ifPDFTeX\else
  % xetex/luatex font selection
  \setmainfont[]{DejaVu Sans}
  \setmonofont[]{DejaVu Sans Mono}
\fi
% Use upquote if available, for straight quotes in verbatim environments
\IfFileExists{upquote.sty}{\usepackage{upquote}}{}
\IfFileExists{microtype.sty}{% use microtype if available
  \usepackage[]{microtype}
  \UseMicrotypeSet[protrusion]{basicmath} % disable protrusion for tt fonts
}{}
\makeatletter
\@ifundefined{KOMAClassName}{% if non-KOMA class
  \IfFileExists{parskip.sty}{%
    \usepackage{parskip}
  }{% else
    \setlength{\parindent}{0pt}
    \setlength{\parskip}{6pt plus 2pt minus 1pt}}
}{% if KOMA class
  \KOMAoptions{parskip=half}}
\makeatother
\usepackage{xcolor}
\usepackage[margin=17mm]{geometry}
\usepackage{color}
\usepackage{fancyvrb}
\newcommand{\VerbBar}{|}
\newcommand{\VERB}{\Verb[commandchars=\\\{\}]}
\DefineVerbatimEnvironment{Highlighting}{Verbatim}{commandchars=\\\{\}}
% Add ',fontsize=\small' for more characters per line
\newenvironment{Shaded}{}{}
\newcommand{\AlertTok}[1]{\textcolor[rgb]{1.00,0.00,0.00}{\textbf{#1}}}
\newcommand{\AnnotationTok}[1]{\textcolor[rgb]{0.38,0.63,0.69}{\textbf{\textit{#1}}}}
\newcommand{\AttributeTok}[1]{\textcolor[rgb]{0.49,0.56,0.16}{#1}}
\newcommand{\BaseNTok}[1]{\textcolor[rgb]{0.25,0.63,0.44}{#1}}
\newcommand{\BuiltInTok}[1]{\textcolor[rgb]{0.00,0.50,0.00}{#1}}
\newcommand{\CharTok}[1]{\textcolor[rgb]{0.25,0.44,0.63}{#1}}
\newcommand{\CommentTok}[1]{\textcolor[rgb]{0.38,0.63,0.69}{\textit{#1}}}
\newcommand{\CommentVarTok}[1]{\textcolor[rgb]{0.38,0.63,0.69}{\textbf{\textit{#1}}}}
\newcommand{\ConstantTok}[1]{\textcolor[rgb]{0.53,0.00,0.00}{#1}}
\newcommand{\ControlFlowTok}[1]{\textcolor[rgb]{0.00,0.44,0.13}{\textbf{#1}}}
\newcommand{\DataTypeTok}[1]{\textcolor[rgb]{0.56,0.13,0.00}{#1}}
\newcommand{\DecValTok}[1]{\textcolor[rgb]{0.25,0.63,0.44}{#1}}
\newcommand{\DocumentationTok}[1]{\textcolor[rgb]{0.73,0.13,0.13}{\textit{#1}}}
\newcommand{\ErrorTok}[1]{\textcolor[rgb]{1.00,0.00,0.00}{\textbf{#1}}}
\newcommand{\ExtensionTok}[1]{#1}
\newcommand{\FloatTok}[1]{\textcolor[rgb]{0.25,0.63,0.44}{#1}}
\newcommand{\FunctionTok}[1]{\textcolor[rgb]{0.02,0.16,0.49}{#1}}
\newcommand{\ImportTok}[1]{\textcolor[rgb]{0.00,0.50,0.00}{\textbf{#1}}}
\newcommand{\InformationTok}[1]{\textcolor[rgb]{0.38,0.63,0.69}{\textbf{\textit{#1}}}}
\newcommand{\KeywordTok}[1]{\textcolor[rgb]{0.00,0.44,0.13}{\textbf{#1}}}
\newcommand{\NormalTok}[1]{#1}
\newcommand{\OperatorTok}[1]{\textcolor[rgb]{0.40,0.40,0.40}{#1}}
\newcommand{\OtherTok}[1]{\textcolor[rgb]{0.00,0.44,0.13}{#1}}
\newcommand{\PreprocessorTok}[1]{\textcolor[rgb]{0.74,0.48,0.00}{#1}}
\newcommand{\RegionMarkerTok}[1]{#1}
\newcommand{\SpecialCharTok}[1]{\textcolor[rgb]{0.25,0.44,0.63}{#1}}
\newcommand{\SpecialStringTok}[1]{\textcolor[rgb]{0.73,0.40,0.53}{#1}}
\newcommand{\StringTok}[1]{\textcolor[rgb]{0.25,0.44,0.63}{#1}}
\newcommand{\VariableTok}[1]{\textcolor[rgb]{0.10,0.09,0.49}{#1}}
\newcommand{\VerbatimStringTok}[1]{\textcolor[rgb]{0.25,0.44,0.63}{#1}}
\newcommand{\WarningTok}[1]{\textcolor[rgb]{0.38,0.63,0.69}{\textbf{\textit{#1}}}}
\usepackage{longtable,booktabs,array}
\usepackage{calc} % for calculating minipage widths
% Correct order of tables after \paragraph or \subparagraph
\usepackage{etoolbox}
\makeatletter
\patchcmd\longtable{\par}{\if@noskipsec\mbox{}\fi\par}{}{}
\makeatother
% Allow footnotes in longtable head/foot
\IfFileExists{footnotehyper.sty}{\usepackage{footnotehyper}}{\usepackage{footnote}}
\makesavenoteenv{longtable}
\setlength{\emergencystretch}{3em} % prevent overfull lines
\providecommand{\tightlist}{%
  \setlength{\itemsep}{0pt}\setlength{\parskip}{0pt}}
\setcounter{secnumdepth}{-\maxdimen} % remove section numbering
\ifLuaTeX
  \usepackage{selnolig}  % disable illegal ligatures
\fi
\usepackage{bookmark}
\IfFileExists{xurl.sty}{\usepackage{xurl}}{} % add URL line breaks if available
\urlstyle{same}
\hypersetup{
  pdftitle={DDTA - Guida integrata per scrittura della documentazione e Base Analysis},
  pdfauthor={Documentation-Driven Threat Analysis - thesis working appendix},
  colorlinks=true,
  linkcolor={Maroon},
  filecolor={Maroon},
  citecolor={Blue},
  urlcolor={Blue},
  pdfcreator={LaTeX via pandoc}}

\title{DDTA - Guida integrata per scrittura della documentazione e Base
Analysis}
\usepackage{etoolbox}
\makeatletter
\providecommand{\subtitle}[1]{% add subtitle to \maketitle
  \apptocmd{\@title}{\par {\large #1 \par}}{}{}
}
\makeatother
\subtitle{Regole, gate, vincoli e protocollo di co-authoring R24}
\author{Documentation-Driven Threat Analysis - thesis working appendix}
\date{Baseline 9cc17a148726bd0734db51e26ac74e031020f340}

\begin{document}
\maketitle

{
\hypersetup{linkcolor=}
\setcounter{tocdepth}{3}
\tableofcontents
}
\textbf{Stato:} CANDIDATE APPENDIX / R24 WORKING INTEGRATION\\
\textbf{Baseline repository di riferimento:}
\texttt{9cc17a148726bd0734db51e26ac74e031020f340}\\
\textbf{Scopo:} consolidare in un unico riferimento stampabile le regole
di authoring della documentazione, i contratti chiusi della Base
Analysis e il protocollo R24 di co-authoring documentazione/BA emerso
dalla sperimentazione BA6.

\begin{quote}
\textbf{Regola di lettura dello stato.} Questa guida distingue sempre
tra: \textbf{CLOSED} (contratto gia chiuso dalla ricerca),
\textbf{R24-WORKING} (regola operativa proposta dall'esperimento
corrente e ancora da validare) e \textbf{CORPUS-SPECIFIC} (vincolo o
convenzione applicabile al caso facial-access, non automaticamente
generalizzabile).
\end{quote}

\subsection{1. Principio di autorita e ciclo documentazione - Base
Analysis}\label{principio-di-autorita-e-ciclo-documentazione---base-analysis}

\subsubsection{1.1 Autorita semantica -
CLOSED}\label{autorita-semantica---closed}

La documentazione governata e l'autorita sul significato del progetto.
La Base Analysis (BA) e una rappresentazione analitica
metodologia-neutrale, ricostruibile e tracciabile di quel significato.

La BA:

\begin{itemize}
\tightlist
\item
  non sostituisce la documentazione;
\item
  non crea automaticamente project truth;
\item
  non completa silenziosamente ambiguita, conflitti o informazioni
  mancanti;
\item
  puo evidenziare lacune, incoerenze o informazioni non rappresentabili;
\item
  puo produrre una proposta di correzione documentale, ma la correzione
  diventa autoritativa solo dopo review e accettazione nella
  documentazione governata.
\end{itemize}

Catena di autorita chiusa da BA3:

\begin{Shaded}
\begin{Highlighting}[]
\NormalTok{governed documentation B0}
\NormalTok{  {-}\textgreater{} accepted BA B0}
\NormalTok{  {-}\textgreater{} downstream analysis / diagnostics}
\NormalTok{  {-}\textgreater{} corrective documentation candidate}
\NormalTok{  {-}\textgreater{} governed review}
\NormalTok{  {-}\textgreater{} governed documentation B1}
\NormalTok{  {-}\textgreater{} accepted BA B1}
\end{Highlighting}
\end{Shaded}

\subsubsection{1.2 Co-authoring come processo di qualita -
R24-WORKING}\label{co-authoring-come-processo-di-qualita---r24-working}

R24 tratta documentazione e BA come \textbf{due processi distinti ma
complementari}:

\begin{Shaded}
\begin{Highlighting}[]
\NormalTok{DOCUMENTATION AUTHORING / REWRITE}
\NormalTok{        |}
\NormalTok{        | authoritative project meaning}
\NormalTok{        v}
\NormalTok{BASE ANALYSIS RECONSTRUCTION}
\NormalTok{        |}
\NormalTok{        | identity / relations / provenance / gaps}
\NormalTok{        v}
\NormalTok{HUMAN{-}READABLE STORY + BA{-}DERIVED VIEW}
\NormalTok{        |}
\NormalTok{        | clarity / completeness / abstraction check}
\NormalTok{        v}
\NormalTok{DOCUMENTATION CORRECTION CANDIDATE}
\NormalTok{        |}
\NormalTok{        \textasciigrave{}{-}{-} governed review {-}{-}\textgreater{} next documentation baseline}
\end{Highlighting}
\end{Shaded}

Il processo e iterativo, ma non circolare sul piano dell'autorita: la BA
puo diagnosticare e proporre; solo la documentazione governata puo
accettare il nuovo significato di progetto.

\subsection{2. Gate preliminare: autorita delle
fonti}\label{gate-preliminare-autorita-delle-fonti}

\subsubsection{2.1 Authority gate - CLOSED /
REQUIRED}\label{authority-gate---closed-required}

Prima di estrarre semantica BA:

\begin{enumerate}
\def\labelenumi{\arabic{enumi}.}
\tightlist
\item
  identificare la baseline immutabile del repository/documentazione;
\item
  identificare i documenti \texttt{CURRENT\_GOVERNED} ammessi come
  primary BA source;
\item
  separare materiale historical, working, non-canonical e regression
  evidence;
\item
  registrare i rami non migrati e gli \texttt{OPEN} noti;
\item
  non usare un documento come primary source solo perche si trova nel
  repository o sembra piu recente.
\end{enumerate}

\textbf{Gate A0 - PASS} solo quando l'autorita di ogni fonte usata e non
ambigua.

\textbf{STOP} se una fonte necessaria non ha stato di autorita
determinabile.

\subsubsection{2.2 Regola historical/current -
CLOSED}\label{regola-historicalcurrent---closed}

Il materiale storico puo essere usato per:

\begin{itemize}
\tightlist
\item
  confronto;
\item
  regression evidence;
\item
  ricostruzione candidata;
\item
  proposta di migrazione.
\end{itemize}

Non puo essere presentato come current project truth finche non e
riscritto, revisionato e accettato nel corpus governato corrente.

\subsection{3. Gerarchia documentale e separazione dei
livelli}\label{gerarchia-documentale-e-separazione-dei-livelli}

Gerarchia della documentazione per il baseline della tesi:

\begin{Shaded}
\begin{Highlighting}[]
\NormalTok{Project problem framing   [precondizione di metodo]}
\NormalTok{        |}
\NormalTok{MacroRequirement}
\NormalTok{        |}
\NormalTok{        +{-}{-} Decision 0..*}
\NormalTok{               |}
\NormalTok{               +{-}{-} FunctionalRequirement 0..*}
\NormalTok{                      |}
\NormalTok{                      +{-}{-} SpecializedRequirement 0..*}
\NormalTok{                             |}
\NormalTok{                             \textasciigrave{}{-}{-} SecurityRequirement [specializzazione corrente]}
\end{Highlighting}
\end{Shaded}

Regole generali:

\begin{itemize}
\tightlist
\item
  la gerarchia e semantica, non dedotta dal nome file o dalla codifica
  dell'ID;
\item
  ogni Decision ha esattamente un owner MR naturale;
\item
  ogni FR ha esattamente una parent Decision;
\item
  ogni Specialized/Security Requirement ha esattamente una parent FR;
\item
  una Decision non e figlia di un FR;
\item
  i livelli non devono essere compressi in un singolo documento o grafo
  se questo distrugge la leggibilita dell'astrazione.
\end{itemize}

\subsection{4. Project problem framing}\label{project-problem-framing}

\textbf{Stato:} CLOSED come precondizione di metodo; non e un nuovo tipo
documentale governato.

Domanda fondamentale:

\begin{quote}
Quale problema generale o classe di problemi deve affrontare il
progetto, entro quale confine significativo, rimuovendo temporaneamente
le soluzioni correnti?
\end{quote}

Regole:

\begin{itemize}
\tightlist
\item
  scrivere una sola volta a livello progetto;
\item
  usare vocabolario di problema/dominio, non architettura corrente;
\item
  non usare provider, device, algoritmo, database, framework o metodo di
  threat analysis come definizione accidentale del problema;
\item
  esplicitare responsabilita importanti fuori scope;
\item
  usare il framing per verificare completezza, sovrapposizione e
  stabilita dell'insieme degli MR;
\item
  non duplicare il framing dentro ogni MR.
\end{itemize}

\textbf{Gate D0 - Project framing}

PASS se un lettore competente nel dominio comprende problema e confine
senza conoscere l'implementazione corrente.

\subsection{5. MacroRequirement - regole di
scrittura}\label{macrorequirement---regole-di-scrittura}

\textbf{Stato:} DOCUMENTATION-LAYER CLOSURE CANDIDATE REVISION 8, usato
come contratto di authoring corrente.

\subsubsection{5.1 Scopo}\label{scopo}

Un MR governa \textbf{una sola responsabilita o concern macro durevole}
necessario ad affrontare il problema di progetto.

Domanda fondamentale:

\begin{quote}
Quale responsabilita/concern macro stabile stiamo governando, quale
risultato o valore deve contribuire, perche e importante e chi
coinvolge?
\end{quote}

\subsubsection{5.2 Semantica richiesta}\label{semantica-richiesta}

\begin{itemize}
\tightlist
\item
  \texttt{id} - identita stabile e tracciabile;
\item
  \texttt{title} - nome conciso comprensibile nel dominio;
\item
  \texttt{lifecycle} - stato current/historical;
\item
  \texttt{intent} - scopo, valore e outcome macro desiderato;
\item
  \texttt{context} - minimo background interpretativo;
\item
  \texttt{stakeholders} - persone/gruppi coinvolti, beneficiari,
  responsabili o impattati;
\item
  \texttt{scope} - confine semantico della responsabilita;
\item
  \texttt{assumptions} - opzionali, solo se valgono per tutto il ramo;
\item
  \texttt{constraints} - opzionali, solo se valgono per tutto il ramo;
\item
  \texttt{dependsOn} - \texttt{MacroRequirement\ {[}0..*{]}} per
  dipendenza/complementarita non gerarchica.
\end{itemize}

\texttt{Objective} resta rimosso: \texttt{Intent} porta scopo, valore e
outcome macro.

\subsubsection{5.3 Invarianti MR}\label{invarianti-mr}

\begin{enumerate}
\def\labelenumi{\arabic{enumi}.}
\tightlist
\item
  nessun parent gerarchico;
\item
  identita indipendente dalla codifica testuale dell'ID;
\item
  ancoraggio al project problem framing;
\item
  un solo macro-concern coerente;
\item
  le Decision figlie restringono il concern, non introducono scope
  estraneo;
\item
  nessun leakage di scelta significativa di soluzione, salvo che sia
  intrinseca al dominio del problema;
\item
  nessuna obbligazione atomica/operativa indipendentemente testabile:
  appartiene a FR;
\item
  leggibilita per stakeholder competenti nel dominio;
\item
  resilienza a cambi di provider/device/algoritmo/deployment;
\item
  stabilita temporale rispetto a cambi di Decision inferiori;
\item
  dipendenze cross-MR esplicite, non pseudo-gerarchie;
\item
  \texttt{title\ +\ intent} dell'intero set MR deve supportare una macro
  project map utile;
\item
  MR abilita l'analisi ma non crea direttamente elementi/metodologie di
  threat analysis;
\item
  una mappa didattica MR non deve essere presentata come DFD,
  architettura o threat model.
\end{enumerate}

\subsubsection{5.4 Split / merge test}\label{split-merge-test}

Prima di dividere o fondere MR, verificare:

\begin{itemize}
\tightlist
\item
  valore indipendente del concern;
\item
  coerenza della futura famiglia di Decision;
\item
  sopravvivenza della distinzione eliminando i nomi della soluzione
  corrente;
\item
  dipendenza vs contenimento;
\item
  spiegabilita da \texttt{title\ +\ intent};
\item
  merge stress: la fusione forzerebbe responsabilita/famiglie di
  Decision sostanzialmente diverse?
\item
  split stress: la divisione nasce da responsabilita di progetto o solo
  da un sottosistema architetturale?
\end{itemize}

\subsubsection{5.5 Gate MR}\label{gate-mr}

\textbf{Gate D1 - MR authoring PASS} quando:

\begin{itemize}
\tightlist
\item
  ogni MR soddisfa gli invarianti sopra;
\item
  l'insieme completo degli MR racconta una storia macro coerente;
\item
  le dipendenze cross-MR necessarie sono esplicite o marcate
  \texttt{OPEN};
\item
  un lettore non ha bisogno delle Decision per capire cosa il progetto
  vuole ottenere a livello macro.
\end{itemize}

\subsection{6. Decision - regole di
scrittura}\label{decision---regole-di-scrittura}

\textbf{Stato:} DOCUMENTATION-LAYER CLOSURE CANDIDATE REVISION 8.

\subsubsection{6.1 Scopo}\label{scopo-1}

Una Decision e un \textbf{commitment significativo governato} che
restringe esattamente un MR fissando una posizione, politica, strategy,
responsibility boundary, convenzione o scelta tecnologico/architetturale
con conseguenze downstream.

Domanda fondamentale:

\begin{quote}
Quale commitment significativo restringe questo MR, perche il progetto
deve prendere questa posizione e quali conseguenze accettiamo?
\end{quote}

\subsubsection{6.2 Core semantico
obbligatorio}\label{core-semantico-obbligatorio}

\begin{Shaded}
\begin{Highlighting}[]
\NormalTok{Decision}
\NormalTok{|{-} title}
\NormalTok{|{-} context}
\NormalTok{|{-} decision}
\NormalTok{\textasciigrave{}{-} consequences}
\end{Highlighting}
\end{Shaded}

Tutti e quattro sono semanticamente richiesti e non null.

\subsubsection{6.3 Regole}\label{regole}

\begin{itemize}
\tightlist
\item
  una Decision = un commitment coerente;
\item
  Context descrive la tensione/alternativa locale ancora aperta prima
  della scelta;
\item
  Decision descrive la posizione scelta, non una tautologia;
\item
  Consequences espone benefici, costi, trade-off, responsabilita e
  obblighi downstream materiali;
\item
  split di commitment che possono cambiare indipendentemente;
\item
  ogni Decision ha un solo MR semantic owner;
\item
  una Decision puo temporaneamente avere \texttt{0..*} FR: zero FR non e
  invalidita automatica, ma e un \textbf{segnale di
  traceability/completeness da revisionare};
\item
  non creare Decision sotto FR.
\end{itemize}

\subsubsection{6.4 Responsibility-boundary
gate}\label{responsibility-boundary-gate}

Prima di scrivere gli FR, chiedere esplicitamente:

\begin{quote}
Il progetto implementa/possiede questa responsabilita oppure consuma una
capability/servizio esterno/organizzativo?
\end{quote}

Il confine cambia quali FR sono responsabilita del progetto.

\subsubsection{6.5 Necessity/default
Decision}\label{necessitydefault-decision}

Consentita solo se, dopo review, non esiste alternativa materialmente
distinta ammissibile.

Deve:

\begin{enumerate}
\def\labelenumi{\arabic{enumi}.}
\tightlist
\item
  spiegare in Context perche le alternative sono assenti/non
  materiali/escluse;
\item
  dichiarare un commitment reale;
\item
  produrre conseguenze downstream materiali;
\item
  essere riapribile se cambia il vincolo o emerge una alternativa;
\item
  non usare ``non abbiamo cercato alternative'' come giustificazione.
\end{enumerate}

\subsubsection{6.6 Gate Decision}\label{gate-decision}

\textbf{Gate D2 - Decision level PASS} quando:

\begin{itemize}
\tightlist
\item
  tutte le Decision rilevanti dei diversi MR sono state esaminate allo
  stesso livello di astrazione;
\item
  ogni Decision restringe chiaramente il proprio MR;
\item
  la storia MR rimane valida e viene solo concretizzata, non riscritta
  accidentalmente;
\item
  i responsibility boundary necessari agli FR sono espliciti;
\item
  le Decision senza FR sono revisionate e motivate oppure registrate
  come gap/candidate.
\end{itemize}

\subsection{7. FunctionalRequirement - regole di
scrittura}\label{functionalrequirement---regole-di-scrittura}

\textbf{Stato:} THESIS WRITING CONTRACT THROUGH FR - REVISION 2.

\subsubsection{7.1 Domanda fondamentale}\label{domanda-fondamentale}

\begin{quote}
Data la parent Decision e una condizione/input applicabile, quale
soggetto/capability deve eseguire quale azione funzionale su/con quale
oggetto/informazione, e quale risultato osservabile o comportamento di
failure deve seguire?
\end{quote}

\subsubsection{7.2 Semantica richiesta}\label{semantica-richiesta-1}

\begin{itemize}
\tightlist
\item
  titolo: un comportamento/capability operativo coerente;
\item
  obbligazione funzionale in prosa normativa leggibile;
\item
  una o piu clausole normative concrete;
\item
  riferimenti strutturati subject/predicate/object come supporto;
\item
  esattamente una parent Decision.
\end{itemize}

\subsubsection{7.3 Regole}\label{regole-1}

\begin{itemize}
\tightlist
\item
  comportamento concreto e indipendentemente valutabile;
\item
  non ripetere la Decision: operationalizzarla;
\item
  preservare condizioni, correlation, lifecycle, atomicity/failure
  semantics nella clausola normativa;
\item
  strutture SPO/relazionali supportano la prosa, non la sostituiscono;
\item
  split di capability che possono cambiare indipendentemente;
\item
  mantenere implementation independence quando piu implementazioni
  valide possono soddisfare lo stesso FR;
\item
  valori mutevoli non vanno hard-coded senza giustificazione semantica;
\item
  riusare riferimenti governati invece di duplicare entita;
\item
  nessun vocabolario di threat method necessario per scrivere FR.
\end{itemize}

\subsubsection{7.4 Cross-MR service
composition}\label{cross-mr-service-composition}

Due MR possono collaborare a un servizio end-to-end senza condividere
ownership degli FR:

\begin{itemize}
\tightlist
\item
  ogni FR rimane in un solo ramo \texttt{MR\ -\textgreater{}\ Decision};
\item
  un consumer FR puo richiedere servizio/capability posseduta da un
  altro MR;
\item
  il consumo non crea un secondo parent;
\item
  se il target canonico del legame non e governato, registrare
  \texttt{OPEN}; non inventarlo nella BA.
\end{itemize}

\subsubsection{7.5 Parameter rule}\label{parameter-rule}

Non introdurre un generico \texttt{functionalParameters}.

Classificare il valore per semantica, per esempio:

\begin{itemize}
\tightlist
\item
  medium/strategy di acquisizione -\textgreater{} Decision;
\item
  latency/throughput -\textgreater{} Performance/Quality candidate;
\item
  retention -\textgreater{} Governance/Privacy candidate;
\item
  provider API limit -\textgreater{} InterfaceContract candidate;
\item
  buffer/thread/tuning -\textgreater{} realization/configuration.
\end{itemize}

\subsubsection{7.6 Gate FR}\label{gate-fr}

\textbf{Gate D3 - FR level PASS} quando:

\begin{itemize}
\tightlist
\item
  ogni FR ha parent Decision unica e coerente;
\item
  ogni commitment operativo necessario e materializzato oppure
  esplicitamente gap;
\item
  i flussi funzionali possono essere raccontati senza ricostruire
  comportamento essenziale da prosa di MR/Decision non
  operationalizzata;
\item
  il modello non richiede di inventare attori, request, store, output o
  bindings non governati.
\end{itemize}

\subsection{8. SpecializedRequirement /
SecurityRequirement}\label{specializedrequirement-securityrequirement}

\textbf{Stato:} SecurityRequirement metamodel S2 CLOSED per il current
thesis baseline.

\subsubsection{8.1 Boundary}\label{boundary}

\texttt{SecurityRequirement\ IS-A\ SpecializedRequirement} e
restringe/enrich un FR senza sostituire la sua funzione.

\subsubsection{8.2 Regole
SecurityRequirement}\label{regole-securityrequirement}

\begin{itemize}
\tightlist
\item
  esattamente una parent FR;
\item
  \texttt{protectedSecurityProperty\ :\ SecurityProperty\ {[}1{]}};
\item
  una sola governing security property per SR;
\item
  clausola normativa con failure mode esplicito;
\item
  cause neutrality: il requisito descrive la proprieta/failure anche se
  la causa puo essere adversarial o non-adversarial;
\item
  non trasformare functional correctness in security solo perche puo
  essere attaccata;
\item
  non usare l'SR come descrizione di attacco;
\item
  non prescrivere automaticamente TLS, mTLS, certificate, VLAN, firma,
  device credential o altri meccanismi: sono realization/Decision
  separate se governate;
\item
  la taxonomy \texttt{SecurityProperty} completa e ancora deferred; non
  inventarla per comodita.
\end{itemize}

\subsubsection{8.3 Gate SR}\label{gate-sr}

\textbf{Gate D4 - Specialized/Security PASS} quando:

\begin{itemize}
\tightlist
\item
  ogni SR ha parent FR unica;
\item
  protegge una sola property governante;
\item
  failure mode e osservabile/esplicito;
\item
  il requisito non confonde funzione, minaccia e realizzazione tecnica;
\item
  l'aggiunta degli SR non altera la storia funzionale, ma la vincola.
\end{itemize}

\subsection{9. Base Analysis - BA1 identity
contract}\label{base-analysis---ba1-identity-contract}

\textbf{Stato:} CLOSED BY BA1-T3.

La BA contiene esattamente due famiglie first-class di identita
semantica:

\begin{Shaded}
\begin{Highlighting}[]
\NormalTok{BAReferent}
\NormalTok{BAProposition}
\end{Highlighting}
\end{Shaded}

\texttt{BAE} e soltanto un termine ombrello, non una terza metaclasse.

\subsubsection{9.1 BAReferent}\label{bareferent}

Unita di significato di progetto metodologia-neutrale, identificabile
indipendentemente e riusabile tra proposition, projection o baseline.

Puo rappresentare participant, capability, information/resource,
behavior/event, contract, store, boundary, state/context ecc. \textbf{La
categoria non crea automaticamente una nuova famiglia BA.}

\subsubsection{9.2 BAProposition}\label{baproposition}

Asserzione metodologia-neutrale identificabile indipendentemente su uno
o piu BAReferent, con identita sufficiente per provenance, diagnosis,
reuse e change disposition.

\subsubsection{9.3 Promotion rule}\label{promotion-rule}

Non creare un BAReferent per ogni nome della prosa.

Un significato locale viene promosso a BAReferent quando necessita
identita/reuse indipendente tra proposition, projection o change. Se una
condizione/valore locale non necessita identita indipendente, puo
restare local value/modifier nei limiti di BA2.

\subsubsection{9.4 Split criterion}\label{split-criterion}

Una nuova famiglia first-class e giustificata solo se entrambe sono
vere:

\begin{enumerate}
\def\labelenumi{\arabic{enumi}.}
\tightlist
\item
  serve identita semantica indipendente;
\item
  invarianti subtype-specific riusabili non possono essere rappresentati
  onestamente con classification, roles, values o proposition sui due
  tipi esistenti.
\end{enumerate}

Convenienza di diagramma, classe software, tabella DB o categoria di un
metodo non basta.

\subsection{10. Base Analysis - BA2 proposition/relation
contract}\label{base-analysis---ba2-propositionrelation-contract}

\textbf{Stato:} CLOSED BY BA2-T4.

\subsubsection{10.1 Shape}\label{shape}

\begin{Shaded}
\begin{Highlighting}[]
\NormalTok{BAProposition}
\NormalTok{|{-} semanticOperatorKey   1}
\NormalTok{|{-} participation         1..*}
\NormalTok{|    |{-} roleKey          1}
\NormalTok{|    \textasciigrave{}{-} term             1}
\NormalTok{|{-} polarity              1}
\NormalTok{\textasciigrave{}{-} scopedModifier        0..* [condition | temporalScope only]}
\end{Highlighting}
\end{Shaded}

\texttt{term} e normalmente un \texttt{BAReferent}; un local value
controllato e ammesso solo se non richiede identita indipendente.

\subsubsection{10.2 Operator registry e role
contract}\label{operator-registry-e-role-contract}

\begin{longtable}[]{@{}
  >{\raggedright\arraybackslash}p{(\columnwidth - 6\tabcolsep) * \real{0.2500}}
  >{\raggedright\arraybackslash}p{(\columnwidth - 6\tabcolsep) * \real{0.2500}}
  >{\raggedright\arraybackslash}p{(\columnwidth - 6\tabcolsep) * \real{0.2500}}
  >{\raggedright\arraybackslash}p{(\columnwidth - 6\tabcolsep) * \real{0.2500}}@{}}
\toprule\noalign{}
\begin{minipage}[b]{\linewidth}\raggedright
Operator
\end{minipage} & \begin{minipage}[b]{\linewidth}\raggedright
Required roles
\end{minipage} & \begin{minipage}[b]{\linewidth}\raggedright
Optional roles
\end{minipage} & \begin{minipage}[b]{\linewidth}\raggedright
Cardinalita
\end{minipage} \\
\midrule\noalign{}
\endhead
\bottomrule\noalign{}
\endlastfoot
\texttt{transfer} & \texttt{source}, \texttt{destination},
\texttt{content} & - & source \texttt{1}; destination \texttt{1..*};
content \texttt{1..*} \\
\texttt{produce} & \texttt{actor}, \texttt{result} & \texttt{input} &
actor \texttt{1}; result \texttt{1..*}; input \texttt{0..*} \\
\texttt{create} & \texttt{actor}, \texttt{created} & - & actor
\texttt{1}; created \texttt{1..*} \\
\texttt{observe} & \texttt{actor}, \texttt{observed} & \texttt{result} &
actor \texttt{1}; observed \texttt{1..*}; result \texttt{0..*} \\
\texttt{transition} & \texttt{subject}, \texttt{toState} &
\texttt{actor}, \texttt{fromState} & subject \texttt{1}; toState
\texttt{1}; actor \texttt{0..1}; fromState \texttt{0..1} \\
\texttt{correlate} & \texttt{correlatedItem},
\texttt{correlationContext} & - & correlatedItem \texttt{1..*};
correlationContext \texttt{1} \\
\texttt{reference} & \texttt{referencer}, \texttt{referenced} & - &
referencer \texttt{1}; referenced \texttt{1..*} \\
\texttt{dependOn} & \texttt{dependent}, \texttt{prerequisite} & - &
dependent \texttt{1}; prerequisite \texttt{1..*} \\
\texttt{consumeService} & \texttt{consumer}, \texttt{service} &
\texttt{provider} & consumer \texttt{1..*}; service \texttt{1}; provider
\texttt{0..1} \\
\texttt{realize} & \texttt{abstract}, \texttt{realization} & - &
abstract \texttt{1}; realization \texttt{1..*} \\
\texttt{assignResponsibility} & \texttt{responsibleParty},
\texttt{responsibilityScope}, \texttt{responsibilityKind} & - & party
\texttt{1..*}; scope \texttt{1}; kind \texttt{1} \\
\texttt{constrain} & \texttt{constraintTarget}, \texttt{constraintValue}
& - & target \texttt{1}; value \texttt{1..*} \\
\texttt{classify} & \texttt{classifiedReferent}, \texttt{semanticKind} &
- & referent \texttt{1}; kind \texttt{1..*} \\
\end{longtable}

\subsubsection{10.3 Polarity}\label{polarity}

Ogni proposition espone:

\begin{Shaded}
\begin{Highlighting}[]
\NormalTok{polarity = affirmative | negative}
\end{Highlighting}
\end{Shaded}

La negazione non crea operatori \texttt{notX}.

\subsubsection{10.4 Scoped modifiers}\label{scoped-modifiers}

Solo due modifier embedded sono ammessi nel lower bound:

\begin{itemize}
\tightlist
\item
  \texttt{condition};
\item
  \texttt{temporalScope}.
\end{itemize}

Possono rimanere embedded solo se sono locali alla proposition, non
introducono participant indipendenti, non richiedono provenance/review
autonoma e non sono riusati altrove.

Failure semantics, atomicity, concurrency, idempotency, correlation,
realization, responsibility, service consumption e reusable constraints
\textbf{non} sono generic modifier: vanno esplicitati con
referent/proposition appropriati.

\subsubsection{10.5 Classification}\label{classification}

\begin{Shaded}
\begin{Highlighting}[]
\NormalTok{classify}
\NormalTok{  classifiedReferent {-}\textgreater{} \textless{}BAReferent\textgreater{}}
\NormalTok{  semanticKind       {-}\textgreater{} \textless{}controlled method{-}neutral key\textgreater{}}
\end{Highlighting}
\end{Shaded}

La taxonomy dei \texttt{semanticKind} non e universalmente chiusa. Il
kind deve avere definizione metodo-neutrale e valore reale per consumer
umani/metodologici, non essere introdotto solo per scegliere un'icona.

\subsubsection{10.6 Lexical separation}\label{lexical-separation}

\begin{Shaded}
\begin{Highlighting}[]
\NormalTok{source wording != stable BA semantic key != display wording}
\end{Highlighting}
\end{Shaded}

\texttt{deliver}, \texttt{send}, \texttt{submit} possono esprimere
\texttt{transfer}; il testo naturale non crea alias normativi.

\subsubsection{10.7 BA2 extension gate}\label{ba2-extension-gate}

Nuovo operator/role solo se un caso governato dimostra che:

\begin{enumerate}
\def\labelenumi{\arabic{enumi}.}
\tightlist
\item
  il significato e metodo-neutrale;
\item
  operatori/roles/polarity/constraints/modifiers esistenti non bastano
  senza perdita;
\item
  la distinzione serve a reuse/consumption/change reasoning;
\item
  non e solo lessicale, implementativa, documentale o method-specific.
\end{enumerate}

Riaprire \textbf{solo} il contratto BA2 minimo interessato.

\subsection{11. Base Analysis - BA3 provenance, review e
change}\label{base-analysis---ba3-provenance-review-e-change}

\textbf{Stato:} CLOSED BY BA3-T4.

\subsubsection{11.1 BAOriginAttachment}\label{baoriginattachment}

Ogni BAReferent/BAProposition materializzato porta o risolve
univocamente a:

\begin{Shaded}
\begin{Highlighting}[]
\NormalTok{BAOriginAttachment}
\NormalTok{  targetElement             1}
\NormalTok{  governedBaselineKey       1}
\NormalTok{  originState               1}
\NormalTok{  sourceLink                0..*}
\NormalTok{  derivationBasisBinding    0..*}
\NormalTok{    inputRoleKey            1}
\NormalTok{    basisRef                1}
\NormalTok{  derivationRuleRef         0..1}
\NormalTok{  revalidationContext       0..*}
\end{Highlighting}
\end{Shaded}

\texttt{sourceLink}, \texttt{derivationBasisBinding} e
\texttt{revalidationContext} hanno responsabilita diverse e non vanno
semanticamente fusi.

\subsubsection{11.2 Origin state}\label{origin-state}

\begin{Shaded}
\begin{Highlighting}[]
\NormalTok{GROUNDED}
\NormalTok{  sourceLink 1..*}
\NormalTok{  no derivation basis/rule}

\NormalTok{DERIVED}
\NormalTok{  derivationBasisBinding 1..*}
\NormalTok{  derivationRuleRef 1}
\NormalTok{  sourceLink optional}

\NormalTok{DIAGNOSTIC\_UNRESOLVED}
\NormalTok{  source/evidence basis \textgreater{}= 1}
\NormalTok{  no silent completion}
\end{Highlighting}
\end{Shaded}

\subsubsection{11.3 Derivation rule}\label{derivation-rule}

Un \texttt{DERIVED} deve puntare a revisione immutabile e ispezionabile
di:

\begin{Shaded}
\begin{Highlighting}[]
\NormalTok{BADerivationRuleDescriptor}
\NormalTok{  ruleKey}
\NormalTok{  ruleRevisionKey [immutable]}
\NormalTok{  inputRoleContract}
\NormalTok{  applicabilityContract}
\NormalTok{  outputSemanticContract}
\NormalTok{  normativeDefinition}
\end{Highlighting}
\end{Shaded}

Logica nascosta di tool/LLM non e sufficiente.

\subsubsection{11.4 Review e freshness}\label{review-e-freshness}

\begin{Shaded}
\begin{Highlighting}[]
\NormalTok{reviewState = PENDING\_REVIEW | ACCEPTED | REJECTED}
\NormalTok{freshness   = CURRENT | STALE}
\end{Highlighting}
\end{Shaded}

Sono dimensioni ortogonali. E consentito, ad esempio,
\texttt{PENDING\_REVIEW\ +\ CURRENT}.

\subsubsection{11.5 Cross-baseline
continuity}\label{cross-baseline-continuity}

\begin{Shaded}
\begin{Highlighting}[]
\NormalTok{RETAIN | REPLACE | RETIRE}
\end{Highlighting}
\end{Shaded}

\begin{itemize}
\tightlist
\item
  BAReferent RETAIN se sopravvive lo stesso significato riusabile
  metodologia-neutrale;
\item
  wording, locator, tecnologia di realization o relazioni mutate non
  forzano da sole nuova identity;
\item
  BAProposition RETAIN solo se operator, polarity, role-bound
  participants/local values e condition/temporalScope restano
  materialmente equivalenti;
\item
  cambi materiali richiedono REPLACE o RETIRE, non editing storico
  silenzioso.
\end{itemize}

\subsubsection{11.6 Change impact}\label{change-impact}

Un change rende candidato \texttt{STALE} solo il significato
esplicitamente dipendente da source, derivation basis, participants,
revalidation context o derivation rule cambiati. Nessuna invalidazione
dell'intera BA per un generico repository change.

\subsection{12. Base Analysis - BA5 canonical
semantics}\label{base-analysis---ba5-canonical-semantics}

\textbf{Stato:} CLOSED BY BA5-T3.

\subsubsection{12.1 Operative semantic
position}\label{operative-semantic-position}

Ogni posizione semanticamente operativa usa token canonico esatto:

\begin{itemize}
\tightlist
\item
  canonical referent name;
\item
  \texttt{semanticOperatorKey};
\item
  \texttt{(semanticOperatorKey,\ roleKey)};
\item
  \texttt{semanticKind};
\item
  controlled value (polarity, modifier kind, responsibility kind, ecc.).
\end{itemize}

La prosa puo restare naturale se non crea un secondo semantic key.

\subsubsection{12.2 Canonical referent
name}\label{canonical-referent-name}

Per baseline governata e naming scope:

\begin{Shaded}
\begin{Highlighting}[]
\NormalTok{one named shared project referent {-}\textgreater{} one exact canonicalName}
\end{Highlighting}
\end{Shaded}

Ma:

\begin{Shaded}
\begin{Highlighting}[]
\NormalTok{BAReferent semantic identity != canonicalName}
\end{Highlighting}
\end{Shaded}

Un rename puo quindi essere \texttt{RETAIN} se il significato e
invariato.

\subsubsection{12.3 No normative synonyms}\label{no-normative-synonyms}

Per il current thesis scope:

\begin{itemize}
\tightlist
\item
  alias/synonym registry normativo: REJECTED AS UNNECESSARY;
\item
  hidden lexical normalization: REJECTED;
\item
  NLP/LLM semantic equivalence come autorita: REJECTED;
\item
  assistenza fuzzy/search/synonym: eventualmente downstream e
  candidate-producing.
\end{itemize}

\subsubsection{12.4 Extension workflow}\label{extension-workflow}

Un token non noto va classificato:

\begin{itemize}
\tightlist
\item
  stesso significato di token esistente -\textgreater{} usare canonical
  token, rifiutare alias;
\item
  method-specific/presentation -\textgreater{} tenere downstream;
\item
  richiede identity riusabile -\textgreater{} BAReferent;
\item
  nuovo semanticKind/controlled value -\textgreater{} evidence +
  definition + review + registry revision immutabile;
\item
  nuovo operator/role semantics -\textgreater{} smallest BA2 reopen.
\end{itemize}

\subsection{13. BA4 / projection
boundary}\label{ba4-projection-boundary}

\textbf{Stato:} boundary chiuso; nessuna nuova sintassi introdotta qui.

\begin{itemize}
\tightlist
\item
  le projection sono downstream della BA;
\item
  una projection puo scegliere layout, aggregazioni e termini
  method-owned;
\item
  non puo creare project truth o modificare BA semantics;
\item
  quando visualizza un shared BAReferent conserva il canonical project
  name;
\item
  una vista didattica, flow view o threat-method view deve essere
  dichiarata come projection, non come fonte BA.
\end{itemize}

\subsection{14. Protocollo R24: analisi breadth-first per
livello}\label{protocollo-r24-analisi-breadth-first-per-livello}

\textbf{Stato:} R24-WORKING - nuova regola operativa da
falsificare/validare.

Il problema osservato in BA6 e che una ricostruzione depth-first
(\texttt{MR\ -\textgreater{}\ Decision\ -\textgreater{}\ FR\ -\textgreater{}\ SR}
su un singolo ramo) puo essere localmente precisa ma distruggere la
storia globale e mescolare livelli di astrazione.

Nuovo ordine:

\begin{Shaded}
\begin{Highlighting}[]
\NormalTok{LEVEL 0  authority + project framing}
\NormalTok{LEVEL 1  ALL MacroRequirements}
\NormalTok{LEVEL 2  ALL Decisions}
\NormalTok{LEVEL 3  ALL FunctionalRequirements}
\NormalTok{LEVEL 4  ALL Specialized/SecurityRequirements}
\NormalTok{LEVEL 5  integrated BA{-}derived views / readiness for overlays}
\end{Highlighting}
\end{Shaded}

\textbf{Regola:} non scendere al livello successivo finche la vista
corrente non e semanticamente coerente e umanamente leggibile, salvo
falsification test esplicitamente isolato.

\subsubsection{14.1 Tre passaggi dentro ogni
livello}\label{tre-passaggi-dentro-ogni-livello}

Per ogni livello:

\textbf{Pass A - Identity / vocabulary}

\begin{itemize}
\tightlist
\item
  identificare significati che richiedono BAReferent;
\item
  riusare identity gia materializzate;
\item
  proporre canonical name senza creare sinonimi;
\item
  classificare solo quando il \texttt{semanticKind} e giustificato
  semanticamente.
\end{itemize}

\textbf{Pass B - Proposition / relations}

\begin{itemize}
\tightlist
\item
  estrarre le relazioni necessarie con BA2;
\item
  non aspettare il grafo finale per accorgersi che i nodi sono
  scollegati;
\item
  non inventare edge per rendere il grafo bello;
\item
  relazioni non esprimibili o non governate diventano diagnostics.
\end{itemize}

\textbf{Pass C - Human story / projection}

\begin{itemize}
\tightlist
\item
  produrre una breve storia in prosa \textbf{solo dai facts BA
  materializzati};
\item
  produrre un grafo/view \textbf{solo dalla BA}, non rileggendo la prosa
  per aggiungere collegamenti mancanti;
\item
  confrontare BA story con source documentation;
\item
  nodi isolati, salti causali e handoff non risolti sono segnali di
  review, non difetti grafici da nascondere.
\end{itemize}

\subsection{15. Gate BA per livello}\label{gate-ba-per-livello}

\subsubsection{15.1 Gate B0 -
authority/provenance}\label{gate-b0---authorityprovenance}

PASS quando ogni elemento BA ha source authority determinabile e
provenance coerente.

\subsubsection{15.2 Gate B1 - MR global
story}\label{gate-b1---mr-global-story}

PASS quando:

\begin{itemize}
\tightlist
\item
  tutti gli MR in scope sono stati analizzati insieme;
\item
  la BA produce una macro story semplice e comprensibile;
\item
  le principali dipendenze/handoff tra concern sono espresse o
  \texttt{DIAGNOSTIC\_UNRESOLVED};
\item
  non servono elementi di Decision/FR per capire il valore macro del
  progetto;
\item
  le views MR non mescolano dettagli architetturali inferiori.
\end{itemize}

\subsubsection{15.3 Gate B2 - Decision
story}\label{gate-b2---decision-story}

PASS quando:

\begin{itemize}
\tightlist
\item
  tutte le Decision in scope raffinano la macro story senza contraddirla
  accidentalmente;
\item
  le principali strategie, responsibility boundary e realization choices
  sono collegate ai concern astratti che concretizzano;
\item
  le Decision indipendenti sono separate;
\item
  i commitment senza conseguenza operativa sono revisionati.
\end{itemize}

\subsubsection{15.4 Gate B3 - Functional
story}\label{gate-b3---functional-story}

PASS quando:

\begin{itemize}
\tightlist
\item
  l'insieme degli FR produce un flow operativo comprensibile;
\item
  request, input, output, actor/capability, correlation e failure
  behavior necessari sono governati o esplicitamente unresolved;
\item
  le cross-MR composition sono traceable senza doppia parentage;
\item
  il grafo funzionale non richiede edge inventati dalla visualizzazione.
\end{itemize}

\subsubsection{15.5 Gate B4 - Specialized/Security
constraints}\label{gate-b4---specializedsecurity-constraints}

PASS quando:

\begin{itemize}
\tightlist
\item
  ogni constraint si attacca a funzione/meaning gia comprensibile;
\item
  security property e failure semantics restano distinti da attack
  description e realization;
\item
  l'aggiunta di security non sostituisce o riscrive la funzione.
\end{itemize}

\subsubsection{15.6 Gate B5 - integrated
readability}\label{gate-b5---integrated-readability}

PASS quando un lettore puo navigare:

\begin{Shaded}
\begin{Highlighting}[]
\NormalTok{macro project story}
\NormalTok{  {-}\textgreater{} decision refinement}
\NormalTok{  {-}\textgreater{} functional flow}
\NormalTok{  {-}\textgreater{} specialized/security constraints}
\end{Highlighting}
\end{Shaded}

senza che un singolo grafo mostri obbligatoriamente tutti i livelli
contemporaneamente.

\subsection{16. Diagnostica: dove appartiene il
problema?}\label{diagnostica-dove-appartiene-il-problema}

\textbf{Stato:} R24-WORKING.

Ogni problema emerso dalla BA deve essere classificato prima di
correggere qualcosa:

\begin{itemize}
\tightlist
\item
  \texttt{DOC\_GAP} - significato necessario ma non
  governato/documentato;
\item
  \texttt{DOC\_STRUCTURE} - informazione presente ma collocata al
  livello/documento sbagliato o con ownership ambigua;
\item
  \texttt{VOCABULARY\_DRIFT} - stesso significato espresso con termini
  incoerenti che impediscono riuso/canonical binding;
\item
  \texttt{BA\_GAP} - documentazione e chiara ma i contratti BA chiusi
  non riescono a preservarne il significato;
\item
  \texttt{PROJECTION\_GAP} - BA e sufficiente ma la view/grafo/story non
  la rende leggibile;
\item
  \texttt{NO\_GAP} - differenza puramente presentazionale o dettaglio
  non governato/non richiesto.
\end{itemize}

\textbf{Regola:} prima si classifica, poi si decide quale layer
riaprire.

\subsection{17. Convenzioni R24 da
validare}\label{convenzioni-r24-da-validare}

\textbf{Stato:} R24-WORKING, non contratti chiusi.

\subsubsection{17.1 Canonical name visibile accanto alla
prosa}\label{canonical-name-visibile-accanto-alla-prosa}

Candidate pattern:

\begin{Shaded}
\begin{Highlighting}[]
\NormalTok{persona presente al punto di accesso \textless{}UserAtGate\textgreater{}}
\NormalTok{identita governata \textless{}GovernedIdentity\textgreater{}}
\end{Highlighting}
\end{Shaded}

Obiettivo: mantenere leggibilita naturale e rendere il riferimento
esatto recuperabile.

Vincoli:

\begin{itemize}
\tightlist
\item
  non introdurre una seconda sintassi semantica;
\item
  il canonical name e governato/registrato solo dopo acceptance;
\item
  durante la ricostruzione puo essere \texttt{PENDING\_REVIEW};
\item
  se la proposta cambia, la documentazione candidate cambia attraverso
  review, non per mutazione invisibile della BA.
\end{itemize}

\subsubsection{17.2 ID basato sulla prima
origine}\label{id-basato-sulla-prima-origine}

Candidate pattern:

\begin{Shaded}
\begin{Highlighting}[]
\NormalTok{MR{-}0003{-}01}
\NormalTok{D{-}3.2{-}03}
\NormalTok{FR{-}3.4.1{-}02}
\end{Highlighting}
\end{Shaded}

Indica il documento in cui l'identity e stata \textbf{prima
materializzata in BA}, non la sua gerarchia semantica. Le occorrenze
successive riusano la stessa identity.

\subsubsection{17.3 Occurrence/sourceLink
index}\label{occurrencesourcelink-index}

Per ogni identity riusata, raccogliere tutte le occorrenze tramite
\texttt{sourceLink}/locator.

Questo puo supportare in futuro indicatori di:

\begin{itemize}
\tightlist
\item
  frequenza;
\item
  dispersione tra documenti/livelli;
\item
  continuita tra baseline;
\item
  stabilita terminologica;
\item
  focus del linguaggio.
\end{itemize}

Tali indicatori sono \textbf{derived diagnostics}, non project semantics
e non criteri automatici di acceptance finche non definiti/validati.

\subsubsection{17.4 semanticKind e
grafica}\label{semantickind-e-grafica}

\texttt{semanticKind} puo essere usato da una projection per scegliere
icone/forme, ma:

\begin{itemize}
\tightlist
\item
  il kind deve essere ammesso per una distinzione metodo-neutrale utile;
\item
  la necessita di un'icona non giustifica un nuovo kind;
\item
  la presentation mapping resta BA4/projection-owned.
\end{itemize}

\subsection{18. Human-readability rules}\label{human-readability-rules}

\textbf{Stato:} R24-WORKING / BA6 pressure.

\begin{enumerate}
\def\labelenumi{\arabic{enumi}.}
\tightlist
\item
  ogni livello deve avere una storia in prosa breve prima di una vista
  dettagliata;
\item
  la prosa di spiegazione e la structured BA devono esprimere lo stesso
  meaning, non diventare due autorita indipendenti;
\item
  una view deve dichiarare cosa include/esclude e quale livello
  rappresenta;
\item
  non mostrare per default MR, Decision, FR e SR nello stesso piano
  grafico;
\item
  il grafo e un test della BA: non deve leggere nuovamente i documenti
  per inventare edge;
\item
  visualizzare chiaramente \texttt{GROUNDED}, \texttt{DERIVED},
  \texttt{DIAGNOSTIC\_UNRESOLVED} quando la distinzione e materialmente
  rilevante;
\item
  \texttt{PENDING\_REVIEW} non significa project truth accettata;
\item
  un nodo isolato puo indicare: identity superflua, relation non
  estratta, documentation gap o abstraction mismatch; va diagnosticato,
  non collegato arbitrariamente;
\item
  la storia end-to-end puo attraversare piu MR: non forzare un singolo
  MR a descrivere responsabilita di altri rami;
\item
  una storia globale incompleta per branch non migrati deve dichiarare
  il confine invece di usare materiale historical come current truth.
\end{enumerate}

\subsection{19. Controlli negativi emersi dal corpus
facial-access}\label{controlli-negativi-emersi-dal-corpus-facial-access}

\textbf{Stato:} CORPUS-SPECIFIC falsification controls, utili come
esempi metodologici.

\subsubsection{19.1 Cross-MR access-decision
consumer}\label{cross-mr-access-decision-consumer}

Il consumer dell'evidenza di verifica appartiene al ramo della decisione
di accesso; se il ramo non e migrato nello stesso baseline:

\begin{itemize}
\tightlist
\item
  mantenere riferimento cross-scope unresolved;
\item
  non inventare canonical target definitivo;
\item
  non usare un componente ipotizzato come project truth.
\end{itemize}

\subsubsection{19.2 Quality/sufficiency
criterion}\label{qualitysufficiency-criterion}

Se il corpus riconosce un criterio di sufficienza/qualita ma non
definisce soglia/requisito specializzato:

\begin{itemize}
\tightlist
\item
  materializzare il gap/candidate;
\item
  non inventare threshold durante BA;
\item
  non nascondere il criterio dentro il significato di un outcome.
\end{itemize}

\subsubsection{19.3 Transfer-to-channel/path
binding}\label{transfer-to-channelpath-binding}

Se \texttt{delivery}, responsibility boundary e medium sono documentati
ma manca il binding strutturale ``THIS transfer uses THIS path'':

\begin{itemize}
\tightlist
\item
  non aggiungere \texttt{channel} a \texttt{transfer} per comodita;
\item
  testare prima se documentazione o derivazione riproducibile possono
  esprimere il binding;
\item
  riaprire BA2 solo con controesempio concreto.
\end{itemize}

\subsubsection{19.4 Intermediate transport
nodes}\label{intermediate-transport-nodes}

Switch, firewall o hop intermedi:

\begin{Shaded}
\begin{Highlighting}[]
\NormalTok{governed + required for declared analysis {-}\textgreater{} must remain recoverable}
\NormalTok{required for declared analysis but undocumented {-}\textgreater{} documentation gap}
\NormalTok{not governed / not required {-}\textgreater{} do not invent}
\end{Highlighting}
\end{Shaded}

\subsubsection{19.5 Test/code evidence
linkage}\label{testcode-evidence-linkage}

Il current baseline non chiude un metamodel di test/code trace.

\begin{itemize}
\tightlist
\item
  non inventare link a test/code nella BA;
\item
  gli FR/SR possono diventare target stabili per future evidence;
\item
  la capacita di collegamento downstream puo essere misurata come
  risultato dell'esperimento, senza creare semantics non governate.
\end{itemize}

\subsection{20. Protocollo di correzione della
documentazione}\label{protocollo-di-correzione-della-documentazione}

Quando la BA trova un problema:

\begin{enumerate}
\def\labelenumi{\arabic{enumi}.}
\tightlist
\item
  registrare il BA diagnostic con evidence/provenance;
\item
  classificare il problema (\texttt{DOC\_GAP}, \texttt{DOC\_STRUCTURE},
  ecc.);
\item
  formulare una \textbf{documentation correction candidate} separata;
\item
  mostrare quale source meaning viene aggiunto/spostato/chiarito;
\item
  non marcare come \texttt{GROUNDED} il nuovo meaning prima
  dell'accettazione documentale;
\item
  sottoporre la correzione a review;
\item
  se accettata, creare nuova governed baseline;
\item
  ricostruire/revalidare BA e applicare BA3
  \texttt{RETAIN\ \textbar{}\ REPLACE\ \textbar{}\ RETIRE}.
\end{enumerate}

\subsection{21. Reopen discipline}\label{reopen-discipline}

Non riaprire un contratto chiuso perche una soluzione sarebbe piu
comoda.

Ordine di diagnosi:

\begin{Shaded}
\begin{Highlighting}[]
\NormalTok{source/documentation insufficiente?}
\NormalTok{  {-}\textgreater{} correggere/proporre documentazione}

\NormalTok{source chiara, BA non esprime il meaning?}
\NormalTok{  {-}\textgreater{} smallest affected BA contract reopen candidate}

\NormalTok{BA sufficiente, view incomprensibile?}
\NormalTok{  {-}\textgreater{} projection/readability correction}
\end{Highlighting}
\end{Shaded}

Nuova family -\textgreater{} BA1 split criterion.\\
Nuovo operator/role -\textgreater{} BA2 extension criterion.\\
Nuova provenance/change need -\textgreater{} BA3 criterion.\\
Problema solo di view -\textgreater{} BA4.\\
Nuovo canonical kind/value -\textgreater{} BA5 evidence-gated
extension.\\
Synonym convenience -\textgreater{} non riaprire BA5.

\subsection{22. Checklist stampabile per una sessione di
lavoro}\label{checklist-stampabile-per-una-sessione-di-lavoro}

\subsubsection{Prima della sessione}\label{prima-della-sessione}

\begin{itemize}
\tightlist
\item[$\square$]
  baseline e authority verificate;
\item[$\square$]
  scope/livello dichiarato;
\item[$\square$]
  tutti i documenti dello stesso livello disponibili;
\item[$\square$]
  historical/current separati;
\item[$\square$]
  OPEN noti caricati.
\end{itemize}

\subsubsection{Durante Pass A -
identity}\label{durante-pass-a---identity}

\begin{itemize}
\tightlist
\item[$\square$]
  nessun noun-to-referent automatico;
\item[$\square$]
  reuse di identity esistenti;
\item[$\square$]
  canonical name unico/candidate;
\item[$\square$]
  semanticKind solo se semanticamente giustificato;
\item[$\square$]
  sourceLink/occurrence registrati.
\end{itemize}

\subsubsection{Durante Pass B -
proposition}\label{durante-pass-b---proposition}

\begin{itemize}
\tightlist
\item[$\square$]
  solo operator BA2 canonici;
\item[$\square$]
  role/cardinality rispettati;
\item[$\square$]
  polarity esplicita;
\item[$\square$]
  local modifier solo condition/temporalScope ammessi;
\item[$\square$]
  nessun edge inventato per il diagramma;
\item[$\square$]
  unresolved -\textgreater{} diagnostic.
\end{itemize}

\subsubsection{Durante Pass C -
story/view}\label{durante-pass-c---storyview}

\begin{itemize}
\tightlist
\item[$\square$]
  storia generata dai facts BA;
\item[$\square$]
  view dello stesso livello di astrazione;
\item[$\square$]
  grounded/derived/unresolved distinguibili;
\item[$\square$]
  nodi isolati revisionati;
\item[$\square$]
  handoff cross-MR espliciti;
\item[$\square$]
  confronto con source documentation completato.
\end{itemize}

\subsubsection{Gate finale del livello}\label{gate-finale-del-livello}

\begin{itemize}
\tightlist
\item[$\square$]
  la storia e comprensibile senza scendere al livello successivo;
\item[$\square$]
  tutte le lacune sono classificate;
\item[$\square$]
  nessuna lacuna e stata silenziosamente riparata;
\item[$\square$]
  documentation candidates separate da accepted BA;
\item[$\square$]
  decisione esplicita: PASS / REWORK / REOPEN CANDIDATE.
\end{itemize}

\subsection{23. Fonti metodologiche
consolidate}\label{fonti-metodologiche-consolidate}

Repository baseline \texttt{9cc17a148726bd0734db51e26ac74e031020f340}:

\begin{itemize}
\tightlist
\item
  \texttt{\_working/ddta-metamodel-working-package/01-mr/01-metamodel/DDTA\_MR\_METAMODEL\_WORKING.md}
\item
  \texttt{02-decision/01-metamodel/DDTA\_DECISION\_METAMODEL\_WORKING.md}
\item
  \texttt{03-functional-requirement/02-authoring-guidance/DDTA\_AUTHORING\_RULES\_THROUGH\_FR\_R2.md}
\item
  \texttt{05-security-requirement/01-metamodel/DDTA\_SECURITY\_REQUIREMENT\_S2\_R1.*}
\item
  \texttt{methodology/BA1\_MINIMAL\_BAE\_IDENTITY\_ONTOLOGY\_R1.md}
\item
  \texttt{methodology/BA2\_RELATION\_ACTION\_VOCABULARY\_R1.md}
\item
  \texttt{methodology/BA3\_PROVENANCE\_DERIVATION\_LIFECYCLE\_CHANGE\_CONTRACT\_R1.md}
\item
  \texttt{methodology/BA5\_CANONICAL\_SEMANTIC\_REGISTRY\_CONTROLLED\_AUTHORING\_CONTRACT\_R1.md}
\item
  \texttt{governed-corpora/facial-access/current/} per i falsification
  controls e la prova R23/R24.
\end{itemize}

\subsection{24. Disposizione R24}\label{disposizione-r24}

Questa guida \textbf{non chiude} R24. Consolida i contratti gia chiusi e
formula il nuovo protocollo di lavoro da testare.

La prossima sperimentazione deve verificare se il co-authoring
breadth-first:

\begin{enumerate}
\def\labelenumi{\arabic{enumi}.}
\tightlist
\item
  preserva autorita e tracciabilita;
\item
  migliora la comprensibilita della storia a ogni livello;
\item
  riduce i nodi BA isolati prodotti da extraction locale;
\item
  distingue meglio lacune documentali da limiti della BA;
\item
  produce una documentazione riscritta piu adatta a review, code/test
  trace futuro e analisi specialistiche downstream.
\end{enumerate}

\end{document}
